\documentclass[a4paper]{article}
\usepackage[pdftex]{graphicx}
\usepackage[utf8]{inputenc}
\usepackage{enumerate}
\usepackage{amssymb}
\usepackage{icomma}
\usepackage{siunitx}
\sisetup{locale=DE} 
\usepackage{href-ul}
\hypersetup{
	colorlinks=true,
	linkcolor=blue,
	urlcolor=blue}
\usepackage{geometry}
\geometry{a4paper, top=15mm, left=15mm, right=15mm, bottom=15mm,
	headsep=10mm, footskip=12mm}

\begin{document}

{\bf Waagrechter Wurf}
\begin{enumerate}[1.]
	\begin{minipage}{0.6\textwidth}
		{\bf \item Aufgabe }
		\begin{enumerate}[a)]
		\item In nebenstehender Anordnung befindet sich eine Kugel in der Position A. A ist doppelt so hoch über dem Boden wie B. 
		Nun wird die Kugel fallengelassen.
		Sie prallt in Punkt B elastisch waagrecht ab und landet schließlich in Punkt C
		auf dem Boden. Berechne die Wurfweite $x$ für $y=\SI{1,0}{\meter}$.
		\item Führe die gleiche Rechnung noch einmal allgemein durch und stelle die Weite $x(y)$ als Funktion von $y$ dar.
		\item Ändert sich die Wurfweite, wenn diese Anordnung auf dem Mond $g_{Mond}=\SI{1,6}{\meter \over {\second^2}}$ betrieben wird?
		Begründe!
		\end{enumerate} 
\end{minipage}
\hspace{0.5 cm}
\begin{minipage}{0.4\textwidth}
	\includegraphics[width=5 cm]{wamm074}		
\end{minipage}

\begin{minipage}{0.6\textwidth}
	{\bf \item Ein Löschflugzeug hat eine Ladung Wasser getankt} und fliegt nun mit einer Geschwindigkeit von $v=\SI{180}{\kilo\meter \over \hour}$ 
	in einer Höhe von $h=\SI{200}{\meter}$ über den Wald. Der Brandherd ist noch klein, er muss genau getroffen werden. In welcher Horizontalentfernung
	muss das Flugzeug die Klappen öffnen, damit das Wasser den Brandherd trifft?
	{\bf \item Sportschütze Lunerald ist sich sicher}, dass er in die Mitte der Zielscheibe gezielt hat. Trotzdem ist der Schuss danebengegangen. Sein Gewehr hat die Kugel auf $v=\SI{310}{\meter\over\second}$ beschleunigt. Wie weit war die Schussdistanz, wenn der zweite Ring der Zielscheibe einen Radius von $\SI{2,0}{\centi\meter}$ hat ?
\end{minipage}
\hfill
\begin{minipage}{0.4\textwidth}
	\includegraphics[width=5 cm]{suss074}		
\end{minipage}

\begin{minipage}{0.6\textwidth}
	{\bf \item Ein Golfball hat einen Durchmesser von} $\SI{43}{\milli\meter}$, der Lochdurchmesser beträgt $\SI{13,4}{\centi\meter}$. Ein Golfball wird dann richtig ,,versenkt", wenn er nicht zu schnell ist. Ansonsten prallt er an der Lochkante nach oben ab.
	\begin{enumerate}[a)]
		\item Welche horizontale und welche vertikale Strecke und fällt er ab der Lochkante, wenn er mit seinem ,,Äquator" am gegenüberliegenden Rand anstoßen soll? 
		\item Berechne hieraus die maximale Geschwindigkeit, die der Ball haben darf, um im Loch zu landen.
	\end{enumerate}  
\end{minipage}
\hspace{0.5 cm}
\begin{minipage}{0.4\textwidth}
	\includegraphics[width=5 cm]{golb074}		
\end{minipage}

\begin{minipage}{0.6\textwidth}
	{\bf \item Spiff, der Klippenspringer nimmt Anlauf...} Wie groß muss seine Horizontalgeschwindigkeit beim Absprung mindestens sein, damit er sicher über die geneigte Klippe ins Wasser fällt? Hierbei sei die Höhe $h=\SI{14}{\meter}$ und er soll zur Sicherheit $\SI{2,0}{\meter}$ neben der Felskante aufkommen.
\vspace{1 cm}

\fbox{
	\begin{minipage}{0.4\textwidth}
		Zur Lösung bitte \href{https://www.okuyakl.de/phys/p9wawL074/ll074.pdf}{hier klicken} oder den QR-Code scannen.\\
	Weitere Arbeitsblätter gibt es unter 
	
	\href{https://www.okuyakl.de}{www.okuyakl.de}
	\end{minipage}
	\hfill
	\begin{minipage}{0.55\textwidth}
		\includegraphics[width=3 cm]{qr074}
		\includegraphics[width=2 cm]{../../viecher/afanticon1}
		
\end{minipage}}

\end{minipage}
\hfill
\begin{minipage}{0.4\textwidth}
	\includegraphics[width=5 cm]{cliff074}		
\end{minipage}

\end{enumerate} 

%------------------------Lösung-------------------------------------------
%\newpage
%\noindent{\bf Aufgabe 1.}\\

\end{document}

